\textbf{(CAP-COLUNI 2018 -- Q11)} A Espanha atual apresenta uma triste realidade na qual a riqueza concentra-se nas mãos de poucos, enquanto muitos encontram-se desempregados. Dados atuais, apontam Amâncio Ortega como o homem mais rico da Espanha, possuindo uma fortuna que ultrapassa a quantia de $x=68\cdot 2^{30}$ dólares, enquanto o alto índice de desemprego atinge cerca de $y=3{,}4$ milhões de espanhóis.

\par\smallskip\textit{(Fonte: Disponíveis em: top10maisricos.org/top-10-mais-ricos-do-mundo/ e https://www.dinheirovivo.pt/. Acesso em: 03 set. 2017. Adaptado).}

Suponha que, em um ato de caridade, Amâncio resolva doar metade da quantia $x$ para ser repartida igualmente entre os $y$ cidadãos espanhóis desempregados. Usando a aproximação $2^{10}\approx10^3$ e sabendo-se que $V$ é a quantia em dólares, doada a cada espanhol desempregado, é CORRETO afirmar a respeito do valor de $V$:

\begin{enumerate}
	\item $950<V<1050$
	\item $1950<V<2050$
	\item $9500<V<10500$
	\item $19500<V<20500$
\end{enumerate}
